\section*{Erd\H{o}s Problem \#366: an ordered pair with different fullness}
\subsection*{1. The exact orientation of the question}
The displayed question asks whether there is a positive integer $n$
such that $n$ is $2$-full and $n+1$ is $3$-full.
The familiar pairs
\[
 (8,9),\qquad(12167,12168)=(23^3,2^3\,3^2\,13^2)
\]
have the opposite orientation: the smaller number is $3$-full,
the larger only $2$-full. They answer an unoriented version but
not the displayed ordered question. This attempt keeps that
distinction, proves an exhaustive parametrization and local sieve,
and establishes conditional finiteness. The existence question
with the displayed orientation remains unresolved.

\subsection*{2. Unique parametrization of $3$-full integers}
A positive integer is $3$-full if and only if it has the unique form
\[
 N=a^3b^4c^5,\quad a,b,c\ge1,\quad
 b,c\text{ squarefree},\quad\gcd(b,c)=1.            \tag{1}
\]
For an exponent $e\ge3$, use $e=3q$ if $e\equiv0\pmod3$,
$e=3q+4$ if $e\equiv1\pmod3$, and $e=3q+5$ if
$e\equiv2\pmod3$. In the latter two cases put the prime in $b$
or $c$ respectively. All resulting $q$ are nonnegative, and
the residue of $e$ makes the choice unique. Primes may also divide $a$.
Thus the number of $3$-full integers up to $X$ is at most
\[
 X^{1/3}\sum_{b\ge1}b^{-4/3}\sum_{c\ge1}c^{-5/3}
 =O(X^{1/3}).                                    \tag{2}
\]
The two convergent series justify this bound without any assumption
on the distribution of prime factors.

Enumerating exactly the triples in (1) up to $10^{10}$ gives $8348$
distinct $3$-full integers. Independently generating the $214122$
$2$-full integers in this range as $u^2v^3$ with $v$ squarefree,
then intersecting the two lists after a shift, gives:
\[
 \begin{array}{c|l}
 \text{condition, with both terms at most }10^{10}
      &\text{smaller terms found}\\ \hline
 n\text{ is }2\text{-full},\ n+1\text{ is }3\text{-full}&\varnothing\\
 n\text{ is }3\text{-full},\ n+1\text{ is }2\text{-full}&8,\ 12167
 \end{array}
\]
This is only a finite certificate. In particular it cannot establish
that the first line is empty without the stated cutoff.

\subsection*{3. Exact local restrictions}
Modulo $p^3$, exclude residues with $v_p(n)=1$ and those with
$v_p(n+1)=1$ or $2$. The respective counts are
$p(p-1)$ and $p(p-1)+(p-1)$. The sets are disjoint, since a prime
cannot divide both consecutive integers. Exactly
\[
 p^3-2p^2+p+1
\]
residues survive this necessary local test. For $p=2$ they are
$n\equiv0,4,7\pmod8$. For any finite set of primes, the Chinese
remainder theorem gives the product of these counts modulo the
product of their cubes. These local conditions are necessary,
not sufficient, since other primes remain untested.

\subsection*{4. Conditional finiteness from $abc$}
Assume that for each $\varepsilon>0$, coprime $a+b=c$ satisfy
$c\le C_\varepsilon\operatorname{rad}(abc)^{1+\varepsilon}$.
For a candidate $n$,
\[
 \operatorname{rad}(n)\le n^{1/2},\qquad
 \operatorname{rad}(n+1)\le(n+1)^{1/3}.
\]
Applying the conjecture with $a=n,b=1,c=n+1$ gives
\[
 n+1\le C_\varepsilon
       \big(n^{1/2}(n+1)^{1/3}\big)^{1+\varepsilon}
 \le C_\varepsilon(n+1)^{5(1+\varepsilon)/6}.
\]
Choose $\varepsilon<1/5$. Then $n$ is bounded, so only finitely many
such oriented pairs can exist under $abc$. The same proof applies
with the orientations exchanged. It does not prove that the
displayed orientation has a single example or has no examples.

\subsection*{5. Assessment and provenance}
Neither the unique parametrization, the finite sieve, nor conditional
finiteness resolves existence. The missing step is an unconditional
construction with the correct orientation, or a uniform exclusion.

Source: \texttt{https://www.erdosproblems.com/366}, checked 2026-09-25;
the page itself flags ambiguity in the original source.
The earlier GPT Pro 5.2 attempt was checked. All unconditional
ingredients here are proved explicitly; $abc$ is an assumption
only in the labeled conditional section.
No novelty or formal Lean verification is claimed.
\subsection*{COMPLETION ESTIMATE}
\textbf{COMPLETION: 35\%}
